% GENERATED FILE. Edit canonical lessons, then run npm run generate:books.
% canonical-source-hash: fnv1a64:3322b0a6e6ea892d
% canonical-lessons: PT-C368-supor, PT-C368-suposicao, PT-C368-hipotese, PT-C368-presumir, PT-C368-plausivel

\chapter{Suponhamos}
\label{ch:pt-suponhamos}
\hlchaptermodality{368}

\begin{chapteropening}
\textbf{By the end of this chapter:} \emph{I can make suppositions and hypotheses, presume something and say whether an idea is plausible.}

Complete the last lesson of chapter 368: I can make suppositions and hypotheses, presume something and say whether an idea is plausible.
\end{chapteropening}

\section[supor]{supor — to suppose}

\label{lesson:PT-C368-supor}



\begin{quote}
Before the new one: say the Portuguese for the voluntary work, then the Portuguese for the charity.
\end{quote}

\subsection*{You'll want to know: supor}
\textbf{supor} — \textquotedblleft{}to suppose\textquotedblright{}.

\begin{quote}
Suponho que ele já saiu de casa.
\end{quote}

\begin{quote}
\emph{I suppose he has already left home.}
\end{quote}

Conjugated like \emph{pôr}: \emph{suponho, supões, supõe}. \emph{Suponhamos que…} — let us suppose that…

\begin{cousinweb}
\emph{supor} is inherited from Latin \emph{suppōnere} \textquotedblleft{}to put or set under, to substitute\textquotedblright{} (EtymDB, Lewis \& Short), built from \emph{sub} and \emph{pōnere} \textquotedblleft{}to put\textquotedblright{}, and reshaped on \emph{pôr}.
\end{cousinweb}

\begin{grammarlens}[title={What you've built}]
One more word for telling, explaining and describing at B1.
\end{grammarlens}

\subsection*{Your turn}
\begin{itemize}
\raggedright
  \item \emph{Say it:} \emph{supor}
  \item \emph{Say it:} \emph{supor}, once more
  \item \emph{Say it:} a sentence of your own with \emph{supor}
  \item \emph{Recall:} say the Portuguese for the voluntary work, then the Portuguese for the charity, then say \emph{supor} again
\end{itemize}

\begin{tcolorbox}[breakable,colback=teal!4,colframe=teal!35!black,title={Before you move on}]
What does supor mean? (\textquotedblleft{}To suppose\textquotedblright{}.) Say it once more.
\end{tcolorbox}
\section[a suposição]{a suposição — the supposition, the assumption}

\label{lesson:PT-C368-suposicao}



\begin{quote}
Before the new one: say the Portuguese for the charity, then the Portuguese for to suppose.
\end{quote}

\subsection*{You'll want to know: a suposição}
\textbf{a suposição} — \textquotedblleft{}the supposition, the assumption\textquotedblright{}.

\begin{quote}
É só uma suposição, não tenho a certeza.
\end{quote}

\begin{quote}
\emph{It is only an assumption; I am not sure.}
\end{quote}

\emph{Fazer suposições} — to make assumptions.

\begin{cousinweb}
\emph{suposição} is a learned borrowing from Latin \emph{suppositiō} \textquotedblleft{}a putting or setting under\textquotedblright{} (EtymDB, Lewis \& Short), from \emph{suppōnere}, on \emph{pōnere} \textquotedblleft{}to put\textquotedblright{}.
\end{cousinweb}

\begin{grammarlens}[title={What you've built}]
One more word for telling, explaining and describing at B1.
\end{grammarlens}

\subsection*{Your turn}
\begin{itemize}
\raggedright
  \item \emph{Say it:} \emph{a suposição}
  \item \emph{Say it:} \emph{a suposição}, once more
  \item \emph{Say it:} a sentence of your own with \emph{a suposição}
  \item \emph{Recall:} say the Portuguese for the charity, then the Portuguese for to suppose, then say \emph{a suposição} again
\end{itemize}

\begin{tcolorbox}[breakable,colback=teal!4,colframe=teal!35!black,title={Before you move on}]
What does a suposição mean? (\textquotedblleft{}The supposition, the assumption\textquotedblright{}.) Say it once more.
\end{tcolorbox}
\section[a hipótese]{a hipótese — the hypothesis; the chance, the possibility}

\label{lesson:PT-C368-hipotese}



\begin{quote}
Before the new one: say the Portuguese for to suppose, then the Portuguese for the supposition.
\end{quote}

\subsection*{You'll want to know: a hipótese}
\textbf{a hipótese} — \textquotedblleft{}the hypothesis; the chance, the possibility\textquotedblright{}.

\begin{quote}
Há uma hipótese de o voo chegar mais cedo.
\end{quote}

\begin{quote}
\emph{There is a chance that the flight will arrive earlier.}
\end{quote}

In Portugal, \emph{hipótese} also means chance: \emph{não tenho hipótese} — I have no chance. \emph{Na hipótese de} — in the event of.

\begin{cousinweb}
\emph{hipótese} is a learned borrowing from Late Latin \emph{hypothesis} (EtymDB), from Greek \emph{hupóthesis} \textquotedblleft{}base, basis of an argument, supposition\textquotedblright{}, built on \emph{hupó} \textquotedblleft{}under\textquotedblright{} and \emph{thésis} \textquotedblleft{}a placing\textquotedblright{}.
\end{cousinweb}

\begin{grammarlens}[title={What you've built}]
One more word for telling, explaining and describing at B1.
\end{grammarlens}

\subsection*{Your turn}
\begin{itemize}
\raggedright
  \item \emph{Say it:} \emph{a hipótese}
  \item \emph{Say it:} \emph{a hipótese}, once more
  \item \emph{Say it:} a sentence of your own with \emph{a hipótese}
  \item \emph{Recall:} say the Portuguese for to suppose, then the Portuguese for the supposition, then say \emph{a hipótese} again
\end{itemize}

\begin{tcolorbox}[breakable,colback=teal!4,colframe=teal!35!black,title={Before you move on}]
What does a hipótese mean? (\textquotedblleft{}The hypothesis; the chance, the possibility\textquotedblright{}.) Say it once more.
\end{tcolorbox}
\section[presumir]{presumir — to presume, to assume}

\label{lesson:PT-C368-presumir}



\begin{quote}
Before the new one: say the Portuguese for the supposition, then the Portuguese for the hypothesis.
\end{quote}

\subsection*{You'll want to know: presumir}
\textbf{presumir} — \textquotedblleft{}to presume, to assume\textquotedblright{}.

\begin{quote}
Presumo que a reunião começa às nove.
\end{quote}

\begin{quote}
\emph{I presume the meeting starts at nine.}
\end{quote}

Formal. \emph{Presumível} — presumed.

\begin{cousinweb}
\emph{presumir} is a learned borrowing from Latin \emph{praesūmere} \textquotedblleft{}to take before, to take beforehand\textquotedblright{} (EtymDB, Lewis \& Short), built from \emph{prae} and \emph{sūmere} \textquotedblleft{}to take\textquotedblright{}.
\end{cousinweb}

\begin{grammarlens}[title={What you've built}]
One more word for telling, explaining and describing at B1.
\end{grammarlens}

\subsection*{Your turn}
\begin{itemize}
\raggedright
  \item \emph{Say it:} \emph{presumir}
  \item \emph{Say it:} \emph{presumir}, once more
  \item \emph{Say it:} a sentence of your own with \emph{presumir}
  \item \emph{Recall:} say the Portuguese for the supposition, then the Portuguese for the hypothesis, then say \emph{presumir} again
\end{itemize}

\begin{tcolorbox}[breakable,colback=teal!4,colframe=teal!35!black,title={Before you move on}]
What does presumir mean? (\textquotedblleft{}To presume, to assume\textquotedblright{}.) Say it once more.
\end{tcolorbox}
\section[plausível]{plausível — plausible, believable}

\label{lesson:PT-C368-plausivel}



\begin{quote}
Before the new one: say the Portuguese for the hypothesis, then the Portuguese for to presume.
\end{quote}

\subsection*{You'll want to know: plausível}
\textbf{plausível} — \textquotedblleft{}plausible, believable\textquotedblright{}.

\begin{quote}
A explicação dele é plausível.
\end{quote}

\begin{quote}
\emph{His explanation is plausible.}
\end{quote}

\emph{Pouco plausível} — hardly believable.

\begin{cousinweb}
\emph{plausível} is a learned borrowing from Latin \emph{plausibilis} \textquotedblleft{}deserving applause, acceptable\textquotedblright{} (EtymDB, Lewis \& Short), built on \emph{plaudere} \textquotedblleft{}to clap\textquotedblright{}: first praiseworthy, then believable.
\end{cousinweb}

\begin{grammarlens}[title={What you've built}]
One more word for telling, explaining and describing at B1.
\end{grammarlens}

\subsection*{Your turn}
\begin{itemize}
\raggedright
  \item \emph{Say it:} \emph{plausível}
  \item \emph{Say it:} \emph{plausível}, once more
  \item \emph{Say it:} a sentence of your own with \emph{plausível}
  \item \emph{Recall:} say the Portuguese for the hypothesis, then the Portuguese for to presume, then say \emph{plausível} again
\end{itemize}

\begin{tcolorbox}[breakable,colback=teal!4,colframe=teal!35!black,title={Before you move on}]
What does plausível mean? (\textquotedblleft{}Plausible, believable\textquotedblright{}.) Say it once more.
\end{tcolorbox}
